% TOP_PROBLEM: {"id": 30006708, "title": "Derandomization of Polynomial Identity Testing", "category_id": 15, "category": {"id": 15, "name": "computer_science", "display_name": "Computer Science", "description": "Computational complexity, algorithms, and theoretical CS.", "slug": "computer-science", "order_index": 15, "created_at": "2026-07-31T15:26:25.671Z"}, "status": "open"}
% FIRST_POSTED: 2026-09-25
% ENTRY_KIND: statement_only
% !TeX program = lualatex
% Definition and sources only; this file is not an LLM solution attempt.
\documentclass[11pt]{article}
\usepackage[margin=25mm]{geometry}
\usepackage{fontspec}
\setmainfont{Latin Modern Roman}
\usepackage{amsmath,amssymb,mathtools}
\usepackage{unicode-math}
\setmathfont{Latin Modern Math}
\usepackage{xurl,hyperref}
\hypersetup{colorlinks=true,urlcolor=blue}
\setcounter{secnumdepth}{0}
\setlength{\parindent}{0pt}
\setlength{\parskip}{5pt}
\setlength{\emergencystretch}{3em}
\begin{document}
\section*{\texorpdfstring{Derandomization of Polynomial Identity Testing}{Derandomization of Polynomial Identity Testing}}\label{30006708}
\subsection{Definitions and mathematical statement}
An arithmetic circuit over a specified effective field has variables and field constants as inputs and addition and multiplication gates. Its output is a polynomial $f_C$. Polynomial identity testing asks whether
\[
 f_C\equiv0\quad\text{in the polynomial ring},
\]
not whether it vanishes at one input or has some root. The goal is a deterministic polynomial-time test in the circuit's encoded size and the necessary field-representation parameters. The catalogue leaves the field, characteristic, and constant-encoding convention implicit; these must be fixed before interpreting ``polynomial time.'' General circuits, not just sparse polynomials or formulas of bounded depth, are intended.
\par\smallskip\textit{Formulation and concept references:} \href{https://arxiv.org/abs/1311.2358}{[S1]}.

\subsection{Short English statement}
Can one deterministically and efficiently recognize when a general arithmetic circuit computes the zero polynomial?
\subsection{Sources}
\begin{itemize}
\item[S1]Derandomizing Polynomial Identity over Finite Fields Implies Super-Polynomial Circuit Lower Bounds for NEXP (Bin Fu).
\url{https://arxiv.org/abs/1311.2358}.
\textit{Formulation source.}
\item[S2]Polynomial Identity Testing for Read-4 Arithmetic Formulas.
\url{https://drops.dagstuhl.de/storage/00lipics/lipics-vol383-ccc2026/LIPIcs.CCC.2026.25/LIPIcs.CCC.2026.25.pdf}.
\textit{Further reference; not a proof endorsement.}
\item[S3]Efficient Polynomial Identity Testing Over Nonassociative Algebras.
\url{https://arxiv.org/abs/2509.11349}.
\textit{Further reference; not a proof endorsement.}
\end{itemize}
\end{document}
